\documentclass{article}
\usepackage{amsmath, amssymb}
\title{Numerical Observations on Radial Scaling of $Q_{i,jk}$}
\date{}
\begin{document}
\maketitle

% Exploratory notes on the radial scaling of the sparse collision operator
% These are numerical findings from n=3 (k=0,1,2); verification at n=4 is pending.

\section*{Numerical Observations on Radial Scaling of $Q_{i,jk}$}

\subsection*{Setup and notation}

The sparse collision operator is stored as a list of $n^3$ matrices:
for each test-function index $i = (k_i, l_i, m_i)$, the matrix
$\bigl[Q_{i,st}\bigr]$ with $s = (k_s, l_s, m_s)$ and $t = (k_t, l_t, m_t)$
gives the bilinear form entries.
The basis functions are $\tilde{f}_{klm} = \mu_{k,l}\,\varphi_{klm}$, where
\[
  \mu_{k,l} = \sqrt{\frac{k!}{\Gamma(k+2l+3)}},
  \qquad
  \varphi_{klm}(r,\theta,\varphi)
    = c_{lm}\,L_k^{2l+2}(r)\,r^l\,P_l^{|m|}(\cos\theta)\,A_m(\varphi),
\]
so every entry satisfies
\[
  Q_{i,st} = \mu_{k_s,l_s}\,\mu_{k_t,l_t}\;\tilde{Q}_{i,st},
\]
where $\tilde{Q}_{i,st}$ is the integral against the un-normalised basis
$\varphi_s$ and $\varphi_t$ (the test function $\varphi_i$ carries no $\mu$
factor in the Petrov--Galerkin formulation).

\subsection*{Structural observations from the $n=3$ operator}

\paragraph{Stable support.}
For a fixed angular pair $(l_i,m_i)$, the set of nonzero positions $(s,t)$
in the operator matrix is \emph{identical} for $k_i = 0,1,2$.
In particular, the nonzero count per test-function slice does not change
with $k_i$.

\paragraph{Sign alternation.}
Every nonzero entry satisfies
\[
  \operatorname{sgn}\bigl(Q_{k_i,st}\bigr) = (-1)^{k_i}\,\operatorname{sgn}\bigl(Q_{0,st}\bigr).
\]
This is consistent with the leading sign of the generalised Laguerre
polynomial $L_{k_i}^{\,2l_i+2}$ through the integrand.

\paragraph{$m$-independence of the scaling.}
Numerically, the ratio $Q_{k_i,st}/Q_{0,st}$ depends only on the six
integers $(k_i,l_i,k_s,l_s,k_t,l_t)$ and is \emph{independent of all
magnetic quantum numbers} $m_i,m_s,m_t$.
This is expected from the rotational structure of the spherical-harmonic
factors, which contribute only Clebsch--Gordan-type constants that cancel
in the ratio.
Verified at $k_s=3$: both $(l_s,l_t)\in\{(1,2),(2,1)\}$ give $S_3/S_0$
independent of $m_s,m_t$ to float64 precision.

\subsection*{Scaling of the stripped integral with radial degree}

Fix the test function at $(k_i = 0, l_i, m_i)$ and one field index, say
$t = (0, l_t, m_t)$, and let $k_s$ vary over $0, 1, 2$.
(The formula below is specific to $k_i = 0$; test functions with $k_i \geq 1$
carry a different $k_s$-scaling pattern that is not characterised here.)
Define the \emph{stripped} sequence
\[
  S_k \;=\; \frac{Q_{i,\,(k,\,l_s,\,m_s),\,t}}
               {\mu_{k,\,l_s}\,\mu_{0,\,l_t}},
\qquad k = 0,1,2,
\]
normalised so that $S_0 > 0$.  (An analogous sequence holds when $k_t$
varies with $k_s$ fixed; the roles of $l_s$ and $l_t$ simply swap.)

For the pairs $(l_s, l_t) \in \{(1,2),(2,1)\}$ the stripped sequence takes
\emph{exact integer} values:

\begin{center}
\begin{tabular}{cccc}
\hline
$l_s$ & $S_0, S_1, S_2$ & $S_3$ (predicted) & $S_3$ (computed) \\
\hline
$1$ & $1,\;-3,\;7$ & $-13$ & $-13$ \\
$2$ & $1,\;-5,\;16$ & $-40$ & $-40$ \\
\hline
\end{tabular}
\end{center}
The $k=3$ values were computed by calling \texttt{operator\_numba} directly
for a single tensor entry (test function $(0,1,0)$, field indices
$(3,l_s,0)$ and $(0,l_t,0)$) using the $(n_\mathrm{lag},n_\mathrm{leb})=(7,9)$
quadrature, without building the full $n=4$ tensor.
After stripping the $\mu$ factors and normalising by $S_0$, both cases
reproduce the formula to float64 precision.

\subsection*{Combinatorial pattern}

With $\alpha_s = 2l_s + 2$ the two sequences above are reproduced by
\[
  \boxed{S_k \;=\; (-1)^k\left[\binom{k+\alpha_s-2}{k}
                              + \binom{k+\alpha_s-4}{k-2}\right]}
\]
where $\binom{n}{k}=0$ whenever $k<0$ or $n<k$.
The first term is $L_k^{\alpha_s-2}(0)$ (the generalised Laguerre polynomial
at the origin); the second is a sub-leading correction that vanishes for
$k<2$.

Sequence values ($k=3$ verified; higher $k$ unverified):
\begin{align*}
  l_s = 1: &\quad S_0,S_1,S_2,\mathbf{S_3},S_4,\ldots
             = 1,-3,7,\mathbf{-13},21,\ldots \\
  l_s = 2: &\quad S_0,S_1,S_2,\mathbf{S_3},S_4,\ldots
             = 1,-5,16,\mathbf{-40},85,\ldots \\
  l_s = 3: &\quad S_0,S_1,S_2,S_3,S_4,\ldots
             = 1,-7,29,-91,238,\ldots
\end{align*}

\subsection*{Extrapolation hypothesis}

If the formula holds for $k\geq 3$, then a tensor entry at radial degree
$k_s = n-1$ (i.e.\ for the next value of $n$) can be recovered without
recomputing the quadrature:
\[
  Q_{i,\,(k_s,\,l_s,\,m_s),\,t}
  \;\approx\;
  S_{k_s}\;\mu_{k_s,\,l_s}\,\mu_{0,\,l_t}
  \;\frac{Q_{i,\,(0,\,l_s,\,m_s),\,t}}{S_0}.
\]
The same applies symmetrically for varying $k_t$.

\subsection*{Connection to the Laguerre recurrence}

The three-term recurrence
\[
  k\,L_k^\alpha(x) = (2k-1+\alpha-x)\,L_{k-1}^\alpha(x)
                   - (k-1+\alpha)\,L_{k-2}^\alpha(x)
\]
evaluated at $x=0$ (where the $x$-term drops out) gives
\[
  k\,A_k = (2k+\alpha_s-3)\,A_{k-1} - (k+\alpha_s-3)\,A_{k-2},
  \qquad \alpha_s = 2l_s+2.
\]
Because $A_k = \binom{k+\alpha_s-2}{k}$ is a binomial coefficient,
this three-term relation degenerates to the simpler \emph{first-order} recurrence
\[
  A_k = \frac{k+\alpha_s-2}{k}\,A_{k-1}, \qquad A_0 = 1,
\]
which generates the stripped sequence $S_k = (-1)^k(A_k + A_{k-2})$ cheaply
without computing factorials or Gamma functions.

\paragraph{Analytic hint from the reduction identity.}
There is a standard Laguerre identity
\[
  L_k^{\alpha-1}(x) = L_k^{\alpha}(x) - L_{k-1}^{\alpha}(x).
\]
Applying it twice at $x=0$ yields
\[
  L_k^{\alpha-2}(0) = L_k^{\alpha}(0) - 2\,L_{k-1}^{\alpha}(0) + L_{k-2}^{\alpha}(0).
\]
The basis functions carry $L_{k_s}^{2l_s+2}$, while the stripped sequence
involves $A_k = L_k^{2l_s}(0)$, a shift of $\alpha$ by $-2$.
Applying the reduction identity twice would produce such a shift with
a two-term remainder, which may explain the structure
$S_k = (-1)^k(A_k + A_{k-2})$.
The precise mechanism — whether it arises from integrations by parts on the
collision kernel or from a generating-function identity — remains to be worked out.

\subsection*{Status and next steps}

\begin{itemize}
  \item $k=0,1,2$ from the stored $n=3$ sparse operator; $k=3$ computed
        directly via a single \texttt{operator\_numba} call (no full $n=4$ build).
        All four data points match the formula to float64 precision.
  \item \textbf{Comprehensive scope check:} a full scan of all 412 entries with
        $k_t=0$, $l_s\geq 1$, $k_s\in\{1,2\}$ in the $n=3$ sparse operator was
        performed.  The 104 entries with $k_i=0$ and
        $(l_s,l_t)\in\{(1,2),(2,1)\}$ all pass (every $m$ combination).
        The remaining 308 entries, which have $k_i\geq 1$ or different
        $(l_s,l_t)$ pairs, follow \emph{different} $k_s$-scaling patterns not
        covered by this formula.
  \item \textbf{Formula scope:} $k_i=0$ is required.  Entries with $k_i=1,2$
        have their own $k_s$-dependence (verified to differ at $k=3$); these
        are not yet characterised.  Similarly, pairs such as
        $(l_s,l_t)\in\{(1,0),(1,1),(2,0),(2,2)\}$ do not follow the same
        formula.
  \item \textbf{Symmetric formula verified:} varying $k_t$ with $k_s=0$ fixed
        yields the same combinatorial formula with $l_t$ replacing $l_s$
        (checked at $k_t=3$: $-13$ for $l_t=1$, $-40$ for $l_t=2$).
  \item The $l_s = 0$ case (isotropic mode) does not follow the same pattern.
  \item The analytic origin of the two-term binomial formula
        $\binom{k+\alpha-2}{k} + \binom{k+\alpha-4}{k-2}$ is unknown;
        it may arise from the Rodrigues representation of $L_k^\alpha$ or
        from a generating-function identity for the collision kernel.
  \item Higher $k$ ($k\geq 4$) and $l_s=3$ remain unverified.
\end{itemize}

\end{document}
